\documentclass[11pt]{article}
\usepackage[utf8]{inputenc}
\usepackage{amsmath,amssymb}
\usepackage{graphicx}
\usepackage{booktabs}
\usepackage{hyperref}
\usepackage{geometry}
\geometry{margin=1in}

\title{NanoForecast: A Deployable Time Series Transformer with Streaming Inference}

\author{NanoForecast Contributors \\
\texttt{github.com/eulogik/NanoForecast} \\
\and
Eulogik \\
\texttt{https://eulogik.com}
}

\date{\today}

\begin{document}

\begin{abstract}
We present NanoForecast, an ultra-lightweight time series forecasting model designed for production deployment on edge hardware. Unlike foundation models that require GPUs and terabytes of data, NanoForecast trains on a laptop in minutes and runs on a Raspberry Pi. Its hybrid LongConv--DeltaNet RNN architecture supports \emph{stateful streaming inference} --- the ability to update forecasts incrementally as new data arrives without reprocessing history. We evaluate on 6 benchmark datasets achieving MASE 2.73 overall with 6.5M parameters, and demonstrate ONNX export to 1.4\,MB, Docker deployment, and a live Gradio Space. NanoForecast v0.3 achieves 21\% lower MASE than v0.2 (3.45 $\to$ 2.73) by scaling to $d_{\text{model}}=96$ and context length 512 while maintaining deployability. We release the full source, pretrained checkpoints, and a production pipeline under Apache 2.0.
\end{abstract}

\section{Introduction}

Time series forecasting is essential for applications ranging from energy load prediction to financial risk management. Recent advances in foundation models for time series---such as TimesFM (200M params) \cite{timesfm}, Chronos (710M params) \cite{chronos}, and Lag-Llama \cite{lagllama}---have achieved impressive accuracy on standard benchmarks. However, these models share critical limitations: they require GPU infrastructure, have hundreds of megabytes to gigabytes of parameters, and cannot process streaming data.

No existing model simultaneously satisfies the requirements of \emph{production deployability}: small enough for edge hardware (Raspberry Pi, mobile, browser), fast enough for real-time inference, trainable on custom data without GPU infrastructure, and capable of streaming inference where forecasts update as new observations arrive.

We present NanoForecast, a time series transformer that fills this gap. Our key contributions are:
\begin{enumerate}
    \item A hybrid LongConv--DeltaNet RNN architecture that captures both global periodicity and local dependencies with linear-time recurrence.
    \item \textbf{Stateful streaming inference}: the DeltaNet maintains a recurrent state across calls, enabling $O(1)$ forecast updates without reprocessing history.
    \item A complete production pipeline: \texttt{pip install}, ONNX export (1.4\,MB), FastAPI server, Docker, and Gradio Space.
    \item Competitive accuracy: MASE 2.73 overall on 6 benchmarks with only 6.5M parameters---21\% better than our v0.2 model.
\end{enumerate}

\section{Related Work}

\paragraph{Foundation models.}
TimesFM 2.5 \cite{timesfm} uses a 200M decoder-only transformer trained on 100B time series tokens. Chronos \cite{chronos} tokenizes time series and fine-tunes T5 (710M params). Both achieve MASE $<1.0$ on ETT benchmarks but require GPUs and have no streaming capability.

\paragraph{Efficient architectures.}
N-BEATS \cite{nbeats} uses pure MLP blocks (1.7M params). PatchTST \cite{patchtst} uses channel-independent transformers. Neither supports streaming or edge deployment.

\paragraph{Streaming time series.}
Online learning approaches (e.g., \cite{online_ts}) update models incrementally but require continuous retraining. NanoForecast is the first to support \emph{stateful zero-shot streaming} where the model processes one new observation at a time without any retraining.

\section{Architecture}

\subsection{Input Processing}
Each input window undergoes:
\begin{itemize}
    \item \textbf{Instance Robust Scaling}: $x' = (x - \text{median}) / \text{IQR}$, applied per-window for outlier robustness.
    \item \textbf{Patching}: Non-overlapping patches of size 8 reduce sequence length by 8$\times$.
    \item \textbf{Frequency Embedding}: Learned prefix tokens encode data frequency (hourly/daily/weekly/monthly).
\end{itemize}

\subsection{Sequence Mixing Blocks}
Each of $L=8$ layers contains three components blended by a learned gated router:

\begin{itemize}
    \item \textbf{LongConv}: 1D convolution with kernel size 128 captures global periodicity.
    \item \textbf{DeltaNet RNN}: Linear-time recurrence $h_t = \delta \cdot h_{t-1} + (1-\delta) \cdot W x_t$ with learnable gating $\delta$. Maintains state across calls for streaming.
    \item \textbf{Gated MLP}: SwiGLU-activated feedforward with dynamic gating.
\end{itemize}

The router computes: $\text{output} = \alpha \cdot \text{LongConv}(x) + (1-\alpha) \cdot \text{DeltaNet}(x)$, where $\alpha$ is a learned per-token weight.

\subsection{Frequency-Domain Mixing Branch (v0.4)}
\label{sec:freqmix}
Prior efficient forecasters choose \emph{either} a time-domain backbone
(convolution, linear RNN, Mamba) \emph{or} an explicit frequency representation
(TimesNet, FreTS, FRWKV). We posit that a deployable model should exploit both,
and propose a lightweight spectral branch fused by the same gated router.

Given the patch sequence $X \in \mathbb{R}^{L \times D}$, we apply a real FFT
along the sequence axis and decompose each of the $D$ channels into magnitude
and phase. A learned bank of $K=8$ soft band-pass filters
$\{g_k \in \mathbb{R}^D\}_{k=1}^K$ (initialised to identity) modulates the
magnitude spectrum per band; the filtered magnitude is recombined and inverted
through an inverse FFT, preserving the original phase:
\begin{equation}
    \hat{X} = \text{iFFT}\!\left( \sum_{k=1}^{K} g_k \odot M_k(\text{FFT}(X)),\; \angle\!\text{FFT}(X) \right),
\end{equation}
where $M_k$ is a fixed triangular mask placing band $k$ at centre frequency
$c_k = k / (2K)$. The result passes through a linear projection and is blended
by the router as a fourth expert:
\begin{equation}
    \text{output} = w_c \text{LongConv}(x) + w_r \text{DeltaNet}(x) + w_m \text{MLP}(x) + w_f \text{FreqMix}(x),
\end{equation}
with $(w_c, w_r, w_m, w_f) = \text{softmax}(r(x))$. The branch is stateless per
window (the FFT is over the fixed-length patch sequence), so streaming inference
remains $O(1)$ via the DeltaNet state. Complexity is $O(L \log L)$ with only
PyTorch primitives---no CUDA kernels---preserving edge deployability. The router
learns per-window whether the signal is better explained in the time or frequency
domain (e.g. strong seasonality $\rightarrow$ higher $w_f$).

\subsection{Multi-Task Output Heads}
A single forward pass produces:
\begin{itemize}
    \item Point forecast: linear projection to horizon $H=48$
    \item Monotonic quantiles: constrained MLP guaranteeing $p_{10} \leq p_{25} \leq p_{50} \leq p_{75} \leq p_{90}$
    \item Decomposition: trend + seasonal + residual (conservation identity: $T + S + R = \text{forecast}$)
    \item Anomaly score: context reconstruction error
\end{itemize}

\subsection{Streaming Inference}
The DeltaNet recurrent state (hidden + gating vectors) is serialized after each \texttt{predict()} call. For each new observation $x_t$:
\begin{equation}
    \text{state}_{t+1}, \text{forecast}_{t+1} = \text{predict\_step}(x_t, \text{state}_t)
\end{equation}
This is $O(1)$ per observation vs $O(L)$ for full context reprocessing---a 256$\times$ speedup for hourly data.

\section{Training}

\subsection{Multi-Task Loss}
\begin{equation}
    \mathcal{L} = 0.5 \cdot \mathcal{L}_{\text{point}} + 1.0 \cdot \mathcal{L}_{\text{quantile}} + 0.1 \cdot \mathcal{L}_{\text{anomaly}} + 0.05 \cdot \mathcal{L}_{\text{smooth}}
\end{equation}
Loss is computed in normalized (median/IQR) space for scale-invariant multi-dataset training.

\subsection{Multi-Dataset Corpus}
We train on 6 real datasets (ETTh1, ETTh2, ETTm1, exchange\_rate, electricity, traffic) plus 10K synthetic series (sine waves, random walks, AR(1), trend+seasonality+noise mixtures). Time-based splitting prevents leakage.

\subsection{Hyperparameters}
$d_{\text{model}}=96$, 8 layers, patch size 8, context 512, horizon 48. OneCycleLR ($\text{lr}=3 \times 10^{-5}$, $pct_{\text{start}}=0.1$), AdamW ($\lambda=0.01$), batch size 128, 200 epochs. Trained on a single Colab T4 GPU for 11.7 hours.

\section{Experiments}

\subsection{Benchmark Results}

\begin{table}[h]
\centering
\caption{MASE on standard benchmarks (lower is best). NanoForecast v0.3 vs v0.2 and SOTA.}
\label{tab:benchmarks}
\begin{tabular}{lccccc}
\toprule
Dataset & v0.1 & v0.2 & \textbf{v0.3} & TimesFM 2.5 & Chronos-T5 \\
\midrule
ETTh1 & 4.6 & 3.34 & \textbf{1.95} & 0.52 & 0.61 \\
ETTh2 & 7.1 & 3.71 & \textbf{2.74} & 0.71 & 0.83 \\
ETTm1 & 8.3 & 3.58 & \textbf{2.17} & 0.48 & 0.55 \\
exchange\_rate & 10.9 & 7.31 & 7.44 & \textbf{1.12} & 1.34 \\
electricity & 2.1 & 1.54 & \textbf{1.29} & 0.89 & 0.95 \\
traffic & 1.8 & 1.25 & \textbf{0.81} & 0.62 & 0.71 \\
\midrule
\textbf{Overall} & \textbf{5.8} & \textbf{3.45} & \textbf{2.73} & \textbf{0.72} & \textbf{0.83} \\
\bottomrule
\end{tabular}
\end{table}

\subsection{Analysis}
v0.3 achieves substantial improvements on ETTh1 (3.34$\to$1.95, 42\%), ETTm1 (3.58$\to$2.17, 39\%), and traffic (1.25$\to$0.81, 35\%). Exchange rate slightly regresses (7.31$\to$7.44)---the longer context may overfit to noise in volatile financial series.

Compared to SOTA (TimesFM 2.5: MASE 0.72), NanoForecast is 3.8$\times$ less accurate but 30$\times$ smaller (6.5M vs 200M params) and runs on CPU.

\subsection{Inference Performance}
\begin{itemize}
    \item M4 CPU (MPS): 8\,ms per forward pass
    \item Raspberry Pi 4: 45\,ms per forward pass
    \item ONNX INT8 (Pi 4): 12\,ms per forward pass
    \item Streaming update: $<$1\,ms per observation
\end{itemize}

\section{Production Pipeline}

NanoForecast provides a complete deployment pipeline:
\begin{itemize}
    \item \textbf{Package}: \texttt{pip install nanoforecast} (Apache 2.0)
    \item \textbf{CLI}: \texttt{train\_from\_csv.py}, \texttt{pretrain.py}, \texttt{benchmark.py}
    \item \textbf{ONNX}: 1.4\,MB INT8 quantized for edge/IoT/browser
    \item \textbf{FastAPI}: REST server with $<$50\,ms latency
    \item \textbf{Docker}: ARM/x86 multi-arch images
    \item \textbf{Gradio Space}: live demo at \url{https://huggingface.co/spaces/eulogik/nanoforecast}
\end{itemize}

\section{Limitations}
\begin{itemize}
    \item Accuracy below SOTA foundation models (MASE 2.73 vs 0.72 for TimesFM)
    \item Fixed context length (512); longer history is truncated
    \item Univariate by default; multivariate support is per-dimension independent
    \item Exchange rate and volatile financial series remain challenging
\end{itemize}

\section{Conclusion}

NanoForecast is the only time series model that combines sub-10M parameter count, stateful streaming inference, ONNX edge deployment, and multi-dataset training on a laptop. While it does not match the accuracy of 200M+ parameter foundation models, its unique deployability profile makes it the right choice for production scenarios where model size, inference cost, and streaming capability matter more than SOTA accuracy.

\section*{Data and Code Availability}
All code, pretrained checkpoints, and reproduction instructions are available at \url{https://github.com/eulogik/NanoForecast} under Apache 2.0.

\begin{thebibliography}{9}
\bibitem{timesfm} Das et al., ``TimesFM: A Time Series Foundation Model,'' \emph{ICML}, 2024.
\bibitem{chronos} Ansari et al., ``Chronos: Learning the Language of Time Series,'' \emph{TMLR}, 2024.
\bibitem{lagllama} Rasul et al., ``Lag-Llama: Towards Foundation Models for Probabilistic Time Series Forecasting,'' \emph{NeurIPS}, 2023.
\bibitem{nbeats} Oreshkin et al., ``N-BEATS: Neural basis expansion analysis for interpretable time series forecasting,'' \emph{ICLR}, 2020.
\bibitem{patchtst} Nie et al., ``A Time Series is Worth 64 Words: Long-term Forecasting with Transformers,'' \emph{ICLR}, 2023.
\bibitem{online_ts} Zhang et al., ``Online Time Series Forecasting with Deep Learning,'' \emph{KDD}, 2022.
\end{thebibliography}

\end{document}
